
\documentclass[11pt,a4paper,oneside]{article}

\synctex=1

%\usepackage[a4,cam,center]{crop}
%\crop[font=\upshape\mdseries\small\textsf]

%% FIXING PROBLEM WITH ALL PAGES PRINTED IN COLOR
%\PassOptionsToPackage{cmyk}{xcolor}% NB: put this *before* \usepackage{pst-all}

\usepackage[utf8]{inputenc} % Para poder escribir con acentos y ñ.
\usepackage[T1]{fontenc}      % Para que haga bien la ``hyphenation''. No
                              % usar si no es necesario, porque ralentiza muchisimo la compilación.
\usepackage{ae}               % Para que todas las fuentes sean Type1, y ninguna Type3.

\usepackage[spanish, english]{babel}

% Use this if you want to include pdf files in the final document
\usepackage[final]{pdfpages}

% Use this if you want to delete headers and footers in empty pages
\usepackage{emptypage}

%\usepackage[nottoc]{tocbibind}
\usepackage{tocbibind}

\usepackage{listings}
\usepackage{longtable}
\usepackage{afterpage}

\usepackage{xspace}
\usepackage{verbatim}
\usepackage{moreverb}
\usepackage{multicol}
\usepackage{amsmath}
\usepackage{eurosym}
\usepackage{subfig} % subfigure is obsolete...
\usepackage{multirow}
\usepackage{fancyhdr}
\usepackage{makeidx}
\usepackage{rotating}
\usepackage{supertabular}
\usepackage{hhline}
\usepackage{array}
\usepackage[noadjust]{cite}     


\usepackage{caption}

% ifthen to allow using language dependent settings
\usepackage{ifthen}

%% JMG: FIXING PROBLEM WITH ALL PAGES PRINTED IN COLOR
% This should not be touched, as it should work as it is know.
\newcommand{\colorspaceused}{rgb}

% %The next section seems to be useless, but it's still pending to try further
% \ifthenelse{\equal{\colorspaceused}{rgb}}
% {
%   \PassOptionsToPackage{rgb}{xcolor}% NB: put this *before* \usepackage{pst-all}
% }
% {
%   \PassOptionsToPackage{cmyk}{xcolor}% NB: put this *before* \usepackage{pst-all}
% }



% To draw rectagles in tfm cover
\usepackage{tikz}


\usepackage{pdftexcmds}

%\usepackage[authoryear]{natbib}
% \makeatletter
% \let\NAT@parse\undefined
% \makeatother
% \usepackage{natbib}

\usepackage{geometry}
\geometry{verbose,a4paper,tmargin=1.5cm,bmargin=1.5cm,lmargin=2cm,rmargin=2cm}
%\geometry{paperwidth=210mm,paperheight=297mm}

\usepackage[
%%ps2pdf,                          %%% hyper-references for ps2pdf
bookmarks=true,%                   %%% generate bookmarks ...
bookmarksnumbered=true,            %%% ... with numbers
hypertexnames=false,               %%% needed for correct links to
                                   %%% figures!!!
%hypertexnames=true,               %%% needed for correct links on pagebackrefs!!!
breaklinks=true,                   %%% breaks lines, but links are very small
%pagebackref=true,
%linktocpage=true,                 %%% enlace en el numero de página.
linktoc=all,
colorlinks=true,
linkcolor=azul,    
citecolor=verde,
urlcolor=azul,                     %%% texto  con color (further
                                   %%% modified in myconfig.tex)
%linkbordercolor={0 0 1},           %%% blue frames around links
pdfborder={0 0 112.0}              %%% border-width of frames 
]{hyperref}                        %%% will be multiplied with 0.009 by ps2pdf


% Para numerar las \subsubsection
\setcounter{secnumdepth}{5}
%para hacer que las \subsubsection aparezcan en el indice
\setcounter{tocdepth}{5}
%\setcounter{lofdepth}{2}
\setcounter{table}{1}
\setcounter{figure}{1}
\setcounter{secnumdepth}{4}

\usepackage{parskip} % Remove first line indent
\usepackage[none]{hyphenat} % avoid hyphenation
%\setlength{\parskip}{1ex plus 0.5ex minus 0.2ex}


\usepackage{multirow}

\usepackage{setspace}
% \renewcommand{\baselinestretch}{10}
\newcommand{\mycaptiontable}[1]{
\begin{spacing}{0.6}
  %\vspace{0.5cm}
  \begin{quote}
    %\begin{center}
            {{Table} \thechapter.\arabic{table}: #1}
    %\end{center}
  \end{quote}
  %\vspace{1cm}
  \end{spacing}
  \stepcounter{table}
}

\newcommand{\mycaptionfigure}[1]{
  %\vspace{0.5cm}
  \begin{spacing}{0.6}
  \begin{quote}
    %\begin{center}
            {{Figure} \thechapter.\arabic{figure}: #1}
    %\end{center}
  \end{quote}
  %\vspace{1cm}
  \end{spacing}
  \stepcounter{figure}
}

\usepackage{amsmath}

\usepackage{courier}

%***************************************************************************
%***************************************************************************
%***************************************************************************
\usepackage{multirow}
\usepackage{rotating}
\usepackage{setspace, amssymb, amsmath, epsfig, multirow, colortbl, tabularx}%
% For acronym package:
% If footnote is specified, text will be included in a footnote
% If printonlyused is specified, only used acronyms will be included
% I use the acronym sty under the sty directory as I needed the newest version
%\usepackage[footnote,printonlyused,withpage]{acronym} 
%\usepackage[printonlyused]{sty/acronym}

% glossaries is better than the acronym package 
%\usepackage[acronym,shortcuts,nomain]{glossaries}

% ifthen to allow using language dependent settings
\usepackage{ifthen}

\newcommand{\clearemptydoublepage}{\newpage{\pagestyle{empty}\cleardoublepage}}

\pagestyle{plain}



\providecommand\phantomsection{}
\onehalfspacing
\sloppy  %better line breaks

%\renewcommand{\chaptermark}[1]{\markboth{\chaptername\ \thechapter.\ #1}{}}
\renewcommand{\sectionmark}[1]{\markright{\thesection\ #1}{}}

%%%%%%%%%%%%%%%%%%%%%%%%%%%%%%%%%%%%%%%%%%%%%%%%%%%%%%%%%%%%%%%%%%%%%%%%%%%
% BEGIN Fancy headers stuff
\fancyhf{}

\fancyhead[LE,RO]{\bfseries\thepage}
\fancyhead[LO]{\bfseries\rightmark}
\fancyhead[RE]{\bfseries\leftmark}

\makeatletter
%\renewcommand{\chaptermark}[1]{\markboth{\@chapapp \ \thechapter . \ #1}{}}
\renewcommand{\sectionmark}[1]{\markright{\thesection \ \ #1}}
\makeatother

\renewcommand{\headrulewidth}{0.5pt}
\renewcommand{\footrulewidth}{0pt}
\addtolength{\headheight}{3.5pt}
\fancypagestyle{plain}{\fancyhead{}\renewcommand{\headrulewidth}{0pt}}
\fancypagestyle{myplain}
{
  \fancyhf{}
  \renewcommand\headrulewidth{0pt}
  \renewcommand\footrulewidth{0pt}
  \fancyfoot[C]{\thepage}
}
% END Fancy headers stuff
%%%%%%%%%%%%%%%%%%%%%%%%%%%%%%%%%%%%%%%%%%%%%%%%%%%%%%%%%%%%%%%%%%%%%%%%%%%

%%%%%%%%%%%%%%%%%%%%%%%%%%%%%%%%%%%%%%%%%%%%%%%%%%%%%%%%%%%%%%%%%%%%%%%%%%%
% BEGIN Set nice chapter titles

% BEGIN Example 0 from http://texblog.org/2012/07/03/fancy-latex-chapter-styles/
% \usepackage[explicit]{titlesec}
% \usepackage{blindtext}
% \definecolor{gray75}{gray}{0.75}
% \newcommand{\hsp}{\hspace{20pt}}
% \titleformat{\chapter}[hang]{\Huge\bfseries}{\chaptername~\thechapter\hsp\textcolor{gray75}{|}\hsp}{0pt}{\Huge\bfseries}
% END Example 0 from http://texblog.org/2012/07/03/fancy-latex-chapter-styles/

% BEGIN Example 1 from http://texblog.org/2012/07/03/fancy-latex-chapter-styles/
% \usepackage{titlesec}
% \usepackage{blindtext}
% \definecolor{gray75}{gray}{0.75}
% \newcommand{\hsp}{\hspace{20pt}}
% \titleformat{\chapter}[hang]{\Huge\bfseries}{\chaptername~\thechapter\hsp\textcolor{gray75}{|}\hsp}{0pt}{\Huge\bfseries}
% END Example 1 from http://texblog.org/2012/07/03/fancy-latex-chapter-styles/

% BEGIN Example 2 from http://texblog.org/2012/07/03/fancy-latex-chapter-styles/
%Options: Sonny, Lenny, Glenn, Conny, Rejne, Bjarne, Bjornstrup
%\usepackage[Sonny]{fncychap}
%\usepackage[Lenny]{fncychap} % ugly
%\usepackage[Glenn]{fncychap}
%\usepackage[Conny]{fncychap} % ugly
%\usepackage[Rejne]{fncychap}
%\usepackage[Bjarne]{fncychap} % Doesn't work in Spanish
%\usepackage[Bjornstrup]{fncychap}
% END   Example 2 from http://texblog.org/2012/07/03/fancy-latex-chapter-styles/

% BEGIN Example 3 from http://texblog.org/2012/07/03/fancy-latex-chapter-styles/
% This is a nice colored example
% \usepackage{kpfonts}
% \usepackage[explicit]{titlesec}
% \newcommand*\chapterlabel{}
% \titleformat{\chapter}
%   {\gdef\chapterlabel{}
%    \normalfont\sffamily\Huge\bfseries\scshape}
%   {\gdef\chapterlabel{\thechapter\ }}{0pt}
%   {\begin{tikzpicture}[remember picture,overlay]
%     \node[yshift=-3cm] at (current page.north west)
%       {\begin{tikzpicture}[remember picture, overlay]
%         \draw[fill=LightSkyBlue] (0,0) rectangle
%           (\paperwidth,3cm);
%         \node[anchor=east,xshift=.9\paperwidth,rectangle,
%               rounded corners=20pt,inner sep=11pt,
%               fill=MidnightBlue]
%               {\color{white}\chapterlabel#1};
%        \end{tikzpicture}
%       };
%    \end{tikzpicture}
%   }
% \titlespacing*{\chapter}{0pt}{50pt}{-60pt}
% END   Example 3 from http://texblog.org/2012/07/03/fancy-latex-chapter-styles/

% BEGIN Example 4 from http://texblog.org/2012/07/03/fancy-latex-chapter-styles/
% END   Example 4 from http://texblog.org/2012/07/03/fancy-latex-chapter-styles/


% END Set nice chapter titles
%%%%%%%%%%%%%%%%%%%%%%%%%%%%%%%%%%%%%%%%%%%%%%%%%%%%%%%%%%%%%%%%%%%%%%%%%%%

%%%%%%%%%%%%%%%%%%%%%%%%%%%%%%%%%%%%%%%%%%%%%%%%%%%%%%%%%%%%%%%%%%%%%%%%%%%
% This is to set background images (in our case to set background image
% in TFMs front and back pages)
% If you want to set this background, use \BgThispage in the
% corresponding pages
% \usepackage[pages=some]{sty/background}
% \backgroundsetup{ scale=1, angle=0, opacity=.1, color=pink,
%   contents={\includegraphics[width=.7\paperwidth]{logos/logoEPS-UAH.jpg}}, vshift=-50pt,  hshift=0pt }


% This is to allow do a clearpage and let the next one to be placed in
% even pages (to set a backpage for example)
\makeatletter
\newcommand*{\cleartoleftpage}{%
  \clearpage
    \if@twoside
    \ifodd\c@page
      \hbox{}\newpage
      \if@twocolumn
        \hbox{}\newpage
      \fi
    \fi
  \fi
}
\makeatother

% Let's define some styles for source code listings:
%
% minimizar fragmentado de listados (from
% http://www.rafalinux.com/?p=599), pero no me funciona:
%\lstnewenvironment{codelisting}[1][]
%{\lstset{#1}\pagebreak[0]}{\pagebreak[0]}
%
% This was using the float package
\usepackage{float}
\floatstyle{plaintop} % optionally change the style of the new float
%\newfloat{codefloat}{H}{cod}[chapter]

\lstdefinestyle{console}
               {
                 basicstyle=\scriptsize\bf\ttfamily,
                 backgroundcolor=\color{gray75},
               }

\lstdefinestyle{Cbluebox}
               {
                 language=C,
                 frame=shadowbox, 
                 rulesepcolor=\color{blue}
               }
                                                
\lstdefinestyle{Cnice}
{
  language=C,
  frame=Ltb,
  framerule=0pt,
  tabsize=2,
  aboveskip=0.5cm,
  framextopmargin=3pt,
  framexbottommargin=3pt,
  framexleftmargin=0.4cm,
  framesep=0pt,
  rulesep=.4pt,
  backgroundcolor=\color{gray97},
  rulesepcolor=\color{black},
  %
  stringstyle=\ttfamily,
  showstringspaces = false,
%  basicstyle=\small\ttfamily,
  basicstyle=\footnotesize\ttfamily,
  commentstyle=\color{gray45},
  keywordstyle=\bfseries,
  %
  numbers=left,
  numbersep=15pt,
  numberstyle=\tiny,
  numberfirstline = false,
  breaklines=true,
}	

\lstdefinestyle{CppExample}
{
  language=C++,
  frame=trbl,
  tabsize=2,
  commentstyle=\textit,
  stringstyle=\ttfamily, 
  basicstyle=\small,
}	

% This one from http://en.wikibooks.org/wiki/LaTeX/Source_Code_Listings
\lstdefinestyle{Ccolor}
{
  belowcaptionskip=1\baselineskip,
  breaklines=true,
  frame=L,
  xleftmargin=\parindent,
  language=C,
  showstringspaces=false,
  basicstyle=\footnotesize\ttfamily,
  keywordstyle=\bfseries\color{green!40!black},
  commentstyle=\itshape\color{purple!40!black},
  identifierstyle=\color{blue},
  stringstyle=\color{orange},
}

% To set side-captions in figures
\usepackage{sidecap}

%%%%%%%%%%%%%%%%%%%%%%%%%%%%%%%%%%%%%%%%%%%%%%%%%%%%%%%%%%%%%%%%%%%%%%%%%%%
% This comes from TeXiS, thanks to its authors, available at
% http://gaia.fdi.ucm.es/projects/texis 
\def\texis{\TeX \raise.15em\hbox{\textsc{i}}S}
%%%%%%%%%%%%%%%%%%%%%%%%%%%%%%%%%%%%%%%%%%%%%%%%%%%%%%%%%%%%%%%%%%%%%%
% Comando:
%
%       \begin{FraseCelebre}
%             \begin{Frase}
%                Y así, del mucho leer y del poco dormir...
%             \end{Frase}
%             \begin{Fuente}
%                Don Quijote de la Mancha
%
%                Miguel de Cervantes
%             \end{Fuente}
%       \begin{FraseCelebre}
%
% Resultado:
%
% Añade la frase célebre del principio de un capítulo.
%%%%%%%%%%%%%%%%%%%%%%%%%%%%%%%%%%%%%%%%%%%%%%%%%%%%%%%%%%%%%%%%%%%%%%
\newenvironment{FraseCelebre}% Definición del entorno de FraseCelebre
 {\begin{list}{}{%
     \setlength{\leftmargin}{0.5\textwidth}% Desplazamos el inicio de
                                 % los párrafos a la derecha la mitad
                                 % de la anchura de la línea de texto.
                                 % Puede que quieras cambiar esto
                                 % por otra cantidad como '5cm'.
     \setlength{\parsep}{0cm}% La separación entre párrafos de la
                             % frase o de la fuente es normal, sin
                             % espacio extra.
     \addtolength{\topsep}{0.5cm}% Aumentamos un poco la separación
                                 % entre la parte de la fase célebre
                                 % y los párrafos de alrededor
   }
 }
 {\unskip \end{list}}

\newenvironment{Frase}%
  {\item \begin{flushright}\small\em}%
  {\end{flushright}}

\newenvironment{Fuente}%
  {\item \begin{flushright}\small}%
  {\end{flushright}}


% To put paragraphs at page bottom
\newenvironment{bottomparagraph}{\par\vspace*{\fill}}{\clearpage}
%\newenvironment{bottomparagraph}{\par\vspace*{\fill}}{\clearemptydoublepage}

\usepackage{wasysym}

% To uppercase first letter
\usepackage{mfirstuc}
% \makefirstuc{abc} produces Abc.
% \makefirstuc{\emph{abc}} produces Abc (\MakeUppercase has been applied to the letter "a" rather than \emph). Note however that
% \makefirstuc{{\em abc}}
% produces ABC (first object is {\em abc} so equivalent to \MakeUppercase{\em abc}), and
% {\makefirstuc{\em abc}}
% produces abc (\em doesn't have an argument therefore first object is \em and so is equivalent to {\MakeUppercase{\em}abc}).
% \makefirstuc{{\'a}bc} produces Ábc.
% \makefirstuc{\ae bc} produces Æbc.
% \makefirstuc{{\ae}bc} produces Æbc.
% \makefirstuc{{ä}bc} produces Äbc. 
%
% Note that non-Latin or accented characters appearing at the start of
% the text must be placed in a group (even if you are using the inputenc
% package). The reason for this restriction is detailed in §3 UTF-8.  
%
% For already defined commands use:
%\newcommand{\abc}{abc}  
%\expandafter\makefirstuc\expandafter{\abc} 

\makeatletter
\def\input@path{{../}{../biblio/}}
%or: \def\input@path{{/path/to/folder}{/path/to/other/folder}}
\makeatother

% This should be relative to the anteproyecto path
\newcommand{\myreferencespath}{../}

%\providecommand{\DIFadd}[1]{{\protect\color{blue}#1}} %DIF PREAMBLE
%\providecommand{\DIFdel}[1]{{\protect\color{red}\protect\scriptsize{#1}}}

% As fancy underlining does not seem to compile with pdflatex, remove underline
%\providecommand{\DIFadd}[1]{{\protect\color{blue}\uwave{#1}}}
\providecommand{\DIFadd}[1]{{\protect\color{blue}\textbf{#1}}}
\providecommand{\DIFdel}[1]{{\protect\color{red}\sout{#1}}}                     
\usepackage{setspace}

% Long underline whole line (http://tex.stackexchange.com/questions/18768/underlining-a-whole-line-with-a-short-text-in-its-beginning)
\usepackage{calc} 
\newlength{\remaining}
\newcommand{\longunderline}[2]{%
\setlength{\remaining}{\linewidth-\widthof{{#2}{#1}}-\tabcolsep}
\noindent\underline{\smash{{#1}\hspace*{\remaining}}}}

\newcommand{\longunderlinea}[2]{%
\setlength{\remaining}{\linewidth-\widthof{{#2}{#1}}-\tabcolsep}
\noindent\underline{\smash{{#1}\hspace*{\remaining}}}}

\newcolumntype{L}[1]{>{\raggedright\let\newline\\\arraybackslash\hspace{0pt}}m{#1}}
\newcolumntype{C}[1]{>{\centering\let\newline\\\arraybackslash\hspace{0pt}}m{#1}}
\newcolumntype{R}[1]{>{\raggedleft\let\newline\\\arraybackslash\hspace{0pt}}m{#1}}

\newcommand{\sixtytextwidth}{0.6\linewidth}
\newcommand{\fortytextwidth}{0.4\linewidth}

%%% Local Variables:
%%% TeX-master: "../book"
%%% End:



